% =============================================================================
%  Lab 12 - CPU Integration
%  THE LONGEST CHAPTER OF THE REPORT. Expect 8-12 pages.
% =============================================================================

\chapter{Lab 12 --- CPU Integration}

\section{Objective}
% Two or three sentences, your own words.

\section{Specification}

\subsection*{The timing contract}
% REQUIRED, in your own words.

\begin{table}[htbp]
\centering
\caption{The two-phase timing contract.}
\label{tab:lab12-timing}
\begin{tabular}{ll}
\toprule
rising edge & \textbf{FETCH} \quad \texttt{IRreg <= ROM[PC]} \\
falling edge & \textbf{COMMIT} \quad register file writes, RAM writes, PC updates \\
\bottomrule
\end{tabular}
\end{table}

% Answer: why does every piece of state commit half a period AFTER the
% instruction is captured? What would break if the register file committed on
% the rising edge instead?

\subsection*{The three wiring decisions}
% State each, and DERIVE the first one rather than quoting it:
%   1. Asel = rt and Bsel = rs. Work through ADDI, SLL, SW and BLEZ and show
%      that one assignment satisfies all of them.
%   2. The store value comes from Aout.
%   3. The ROM is addressed by PCreg, not NSA.

\section{Design}
% Answer all three:
%   (a) Why do the register file and the NIAG need an INVERTED clock? What does
%       inverting at the instantiation let you avoid doing?
%   (b) If Asel and Bsel were swapped, three of the four instruction types would
%       still work. Which one breaks, and why does that make the bug harder to
%       find rather than easier?
%   (c) Your CPU has ONE functional pin. In Lab 5 the ALU needed 108 as a
%       top-level entity. Explain the difference.

\section{Implementation}

\vhdlfile{code/lab12/CPU.vhd}{\texttt{CPU} --- twelve components, wired}

% State clearly which components are your own and which are TA reference
% versions. Using a reference component costs nothing - but the report must be
% honest about what you built.

\section{Verification}

\subsection*{Hand-computed trace}
% REQUIRED, and done BEFORE simulating. Derive these from the ROM listing and
% your control table, then put the simulated values beside them.

\begin{table}[htbp]
\centering
\caption{The first five cycles: predicted by hand, then measured.}
\label{tab:lab12-hand}
\begin{tabular}{cllllll}
\toprule
cyc & PC & instruction & ALU (hand) & ALU (sim) & reg written & next PC \\
\midrule
0 & \texttt{0x00} & \texttt{LUI r1, 11} & & & & \\
1 & \texttt{0x04} & & & & & \\
2 & \texttt{0x08} & & & & & \\
3 & \texttt{0x0C} & & & & & \\
4 & \texttt{0x10} & & & & & \\
\bottomrule
\end{tabular}
\end{table}

\subsection*{The full execution trace}
% Paste the complete 45-row trace the testbench prints. This is the evidence
% that the machine runs.

\subsection*{Final state}

\begin{table}[htbp]
\centering
\caption{Final state after 45 instructions.}
\label{tab:lab12-final}
\begin{tabular}{lll}
\toprule
quantity & expected & measured \\
\midrule
instructions executed & 45 & \\
CPI & 1 & \\
final \texttt{r1} & \texttt{0x00000000} & \\
final \texttt{r2} & \texttt{0x00000400} & \\
word stored at address 0 & \texttt{0x00000400} & \\
\bottomrule
\end{tabular}
\end{table}

% ONE SENTENCE: what number is 0x400, and where does it come from?

\section{Results}
% ALL SEVEN evidence sources. This is the only design in the course that
% reports every one of them meaningfully. One sentence of interpretation each.

\figplaceholder{the annotated waveform, about 20 cycles, with the instruction
boundaries marked and at least one loop iteration visible end to end}

\figplaceholder{the RTL Viewer, showing the twelve components and how they connect}

% E1: errors and warnings. Explain the metavalue warnings: why they appear in
%     the first few cycles and why they stop.
% E3: Technology Map - what the fabric actually made of it.
% E4: registers (PC 32, IR 64, register file ~1024) and block memory bits.
% E5: the pin count. ONE functional pin. Compare with the ALU's 108 in Lab 5
%     and say what that comparison means.
% E6: Fmax on the Slow 1100mV 85C model, and the critical path. Which component
%     is on it, and is that what you predicted?

\section{Discussion}
% REQUIRED, and the most important Discussion of the semester.
%
% For EACH integration bug you hit:
%   - what the symptom was
%   - WHICH OF THE FOUR CHECKS found it
%   - what you would look for first if you saw that symptom again
%
% Nobody wires twelve components correctly first time. A Discussion claiming
% otherwise scores zero.
%
% Then close with one paragraph: what did you actually learn from this lab that
% you did not learn from building the components?
